% ===================================================================
%  BAGAN No. 9 — Gunung. Hutan.
%
%  Traduction, faite en 2026, en indonesien d'aujourd'hui, sur l'ido
%  de texto/io. Regles de la colonne : voir 10-bagan-01.tex. Deux
%  scenes, deux numerotations : les cles portent c1 et c2. Le bloc
%  t09-c1-05-1 tient DEUX alineas sous une seule cle — ordo.py le
%  marque « ¶2 », apres deux cles inventees dans la colonne ourdoue.
%
%  LE VOLCAN EST INDONESIEN JUSQU'AU BOUT, ET LES CONIFERES NE LE SONT
%  PAS. C'est le partage le plus net de tout le tableau : gunung
%  berapi la montagne de feu (6), kawah le cratere, lereng la pente
%  (7), jurang le ravin (32), air terjun la cascade (10), lava — tout
%  ce qui touche au volcan porte un nom malais transparent, dans un
%  pays qui compte plus de volcans actifs que tout autre. Mais le
%  sapin (11, 31), le cypres, le tilleul, le hetre, le bouleau,
%  l'orme, le cedre s'empruntent tous : ce sont des arbres de la zone
%  temperee, et la montagne de Java n'est pas faite du meme bois que
%  la montagne des Pyrenees. Une meme planche, deux resultats
%  opposes, et la ligne passe entre la GEOLOGIE et la BOTANIQUE.
%
%  ET gletser (5) CONFIRME L'OURDOU MOT POUR MOT. La colonne ourdoue
%  avait donne le contre-exemple du releve : گلیشیئر est un emprunt
%  dans la langue d'un pays qui porte le Baltoro. L'Indonesie a elle
%  aussi des glaciers — ceux du Puncak Jaya, les seuls glaciers
%  tropicaux d'Asie — et elle dit gletser, du neerlandais. Deux
%  langues, deux pays qui ont la chose, deux emprunts. La regle posee
%  au tableau 10 ourdou tient : ce n'est pas la presence de la chose
%  qui decide du mot, c'est la DATE a laquelle on a eu besoin de la
%  nommer par ecrit.
%
%  damar (27) VA DANS L'AUTRE SENS, ET C'EST LE QUATRIEME MOT DU
%  RELEVE A LE FAIRE. La resine des coniferes est damar en malais, et
%  l'anglais en a fait « dammar », terme de peinture et de vernis
%  employe dans le monde entier. Apres برآمدہ la veranda, قلی le
%  porteur et بازار le bazar de la colonne ourdoue, c'est le premier
%  mot que l'indonesien donne au vocabulaire europeen — et il le donne
%  precisement sur le seul objet de cette foret qui soit aussi un
%  produit d'exportation.
%
%  murai (56), LA PIE : LA OU L'OURDOU AVAIT DU EMPRUNTER. La colonne
%  ourdoue ecrivait میگپائی faute de mot descendu dans la plaine ;
%  l'indonesien a murai, oiseau noir et blanc de l'archipel, connu de
%  tous. Avec burung pelatuk le pic (30), tupai l'ecureuil (61),
%  paku-pakuan les fougeres (55), la foret garde ses noms de fond
%  partout ou l'espece existe — et n'emprunte que les arbres.
%
%  pal (bloc t09-c1-15-1) DEMANDE UNE NOTE. L'ido dit « quaropa
%  kilometri », la lieue de quatre kilometres. L'indonesien a le pal,
%  unite des routes des Indes neerlandaises valant 1 507 metres, qui
%  survit dans des noms de lieux comme Palmerah. Trois pals font une
%  lieue a peu de chose pres. Troisieme colonne du releve a mesurer
%  cette unite, apres le காதம் tamoul (dix a seize kilometres) et le
%  کوس ourdou (trois) : la lieue ne se traduit jamais deux fois de la
%  meme facon.
% ===================================================================

%%K t09-apar-1 apar
\VUsaut{4.0mm}
\VUtitre{15.0pt}{60}{BAGAN No. 9}
\VUsaut{3.0mm}

%%K t09-tit-1 sub
\VUcentre{12.0pt}{0}{\textit{Adegan Pertama.}}
\VUsaut{3.0mm}

%%K t09-tit-2 sub
\VUpk{11.4pt}{40}{Gunung. --- Kota pemandian.}
\VUsaut{3.0mm}

%%K t09-c1-01-1 p
1. --- Sudah delapan hari saya berada di Pratz. Saya tidak dapat
mengungkapkan seluruh kekaguman yang saya rasakan ketika saya berdiri
di hadapan \VUgras{pegunungan} \textsuperscript{(1)} Pyrenea yang
menakjubkan, yang \VUgras{punggungnya} \textsuperscript{(2)} tampak di
langit biru seperti untaian raksasa.

%%K t09-c1-02-1 p
2. --- Saya tidak membayangkan bahwa \VUgras{puncak-puncak}
\textsuperscript{(3)} Pyrenea mencapai \VUgras{ketinggian} yang begitu
besar, dan saya tercengang di hadapan \VUgras{gletser-gletser}
\textsuperscript{(5)} yang megah dan \VUgras{puncak-puncak runcing}
\textsuperscript{(4)} bersalju, tempat \VUgras{longsoran salju} yang
begitu mengerikan berguling turun.

%%K t09-tit-3 sub
\VUsaut{3.0mm}
\VUpk{10.2pt}{40}{Pemandangan gunung.}
\VUsaut{3.0mm}

%%K t09-c1-03-1 p
3. --- Sejak hari sesudah kedatangan saya, saya pergi ke
\VUgras{gunung berapi} \textsuperscript{(6)} tua yang sudah padam di
dekat Pratz. Di tepi \VUgras{kawah} yang gersang dan gundul, orang
masih jelas melihat jejak jalannya \VUgras{lava} yang mengalir,
beberapa abad yang lalu, di \VUgras{lereng gunung}
\textsuperscript{(7)}. Saya berhasil menemukan, di salah satu
\VUgras{lerengnya,} beberapa pecahan lava itu.

%%K t09-c1-04-1 p
4. --- Pratz terletak di perbatasan, dan \VUgras{benteng}
\textsuperscript{(8)} yang menjaga \VUgras{celah gunung}
\textsuperscript{(27)} yang memisahkan Prancis dari Spanyol sungguh
mengingatkan pada kastil-kastil tua di tepi Rhein, yang tampak
mengintai \VUgras{mangsa} dari puncak batu karangnya.

%%K t09-c1-05-1 p
5. --- Seekor \VUgras{elang} \textsuperscript{(9)} yang sangat besar,
yang sarangnya tidak jauh dari \VUgras{benteng} \textsuperscript{(8)},
sering menarik perhatian saya.

\VUgras{Burung pemangsa} itu, yang \VUgras{paruhnya} kuat, sering
membawa untuk anak-anaknya seekor anak domba atau anak kambing yang
dicengkeramnya dengan \VUgras{cakarnya.} Begitu besar \VUgras{sayapnya}
sehingga ia dapat lama melayang dengan bebannya yang berat.

%%K t09-tit-4 sub
\VUsaut{3.0mm}
\VUpk{10.2pt}{40}{Pelancong.}
\VUsaut{3.0mm}

%%K t09-c1-06-1 p
6. --- Di sana tamasya yang paling menyenangkan adalah perjalanan ke
air terjun \textsuperscript{(10)} « Kau ». \VUgras{Curahan air} itu
jatuh berpusar lalu lenyap sebagai \VUgras{anak sungai} yang indah di
\VUgras{hutan pohon fir} \textsuperscript{(11)} di sebelah barat kota.
Kami mendaki lewat hutan itu sampai ke \VUgras{observatorium cuaca}
\textsuperscript{(12)}, dan dari puncak \VUgras{menara pandang} kami
melihat sekilas seluruh \VUgras{pemandangan luas} daerah itu.
\VUgras{Pemandangannya} menakjubkan, dan \VUgras{alam} tampak
mencurahkan kepada negeri itu segala harta yang dimilikinya.

%%K t09-c1-07-1 p
7. --- Untuk turun, kami memakai \VUgras{kereta gantung}
\textsuperscript{(13)} dan berhenti di sebuah \VUgras{stasiun} kecil di
dekat \VUgras{reruntuhan} \textsuperscript{(14)} sebuah \VUgras{kastil
feodal} tua yang dahulu ditempati para tuan Pratz.

%%K t09-c1-08-1 p
8. --- Kasihan kastil itu! Apa gerangan yang dipikirkannya ketika
melihat di bawah sana \VUgras{kota pemandian} \textsuperscript{(15)}
yang genit, yang kini menjadi \VUgras{tempat mandi air panas} yang
sedang digemari?

%%K t09-c1-08-2 p
\VUgras{Iklimnya} begitu menyenangkan dan \VUgras{air mineralnya}
begitu manjur sehingga \VUgras{pengobatan} yang dijalani di
\VUgras{sanatorium} \textsuperscript{(17)} sungguh menjadi kesenangan.

%%K t09-tit-5 sub
\VUsaut{3.0mm}
\VUpk{10.2pt}{40}{Kota air panas yang menyenangkan.}
\VUsaut{3.0mm}

%%K t09-c1-09-1 p
9. --- Lagi pula, segala sesuatu dikerjakan untuk membuat masa tinggal
menyenangkan. Sebuah \VUgras{jembatan layang} \textsuperscript{(16)}
yang besar dibangun untuk memungkinkan jalan kereta api ;
\VUgras{taman} \textsuperscript{(18)} \VUgras{gedung pemandian,} tempat
orang mengadakan \VUgras{konser,} ditata dengan cara yang memesona.
Beberapa « \VUgras{vila} » \textsuperscript{(19)} yang anggun dibangun
di sekeliling kasino \textsuperscript{(21)}, dan sebuah \VUgras{jalan
raya} \textsuperscript{(22)} yang terpelihara sangat baik amat
memudahkan perhubungan.

%%K t09-c1-10-1 p
10. --- Sebuah \VUgras{tanggul} \textsuperscript{(23)} sederhana
memisahkan \VUgras{jalan} \textsuperscript{(20)} dari \VUgras{lembah}
\textsuperscript{(25)} yang sebenarnya, yang di dasarnya terletak
sebuah \VUgras{danau} \textsuperscript{(26)} kecil yang anggun, yang
memberi air kepada sebuah \VUgras{anak sungai} \textsuperscript{(24)}
yang jernih.

%%K t09-c1-11-1 p
11. --- Di sisi lembah yang lain, sebuah \VUgras{tambang besi}
\textsuperscript{(28)} digarap. Beberapa kali saya pergi melihat
\VUgras{para penambang} turun ke \VUgras{sumur tambang,} dan saya
bahkan masuk ke \VUgras{terowongan} tempat mereka menghabiskan hari
mereka, menggali \VUgras{bijih} yang mengandung \VUgras{logam.}

%%K t09-c1-12-1 p
12. --- Sebuah \VUgras{peleburan} besar dibangun di dekat tambang itu
untuk segera memanfaatkan kekayaan yang digali darinya. \VUgras{Tanur
tinggi} \textsuperscript{(29)}, yang dipanaskan dengan \VUgras{batu
bara} yang melimpah di sini, memungkinkan memperoleh \VUgras{besi
tuang} dengan cepat dan murah.

%%K t09-c1-13-1 p
13. --- Tidak jauh dari tambang itu orang menggali sebuah
\VUgras{tambang batu} \textsuperscript{(30)}. Bukan \VUgras{granit,}
bukan \VUgras{porfir,} bukan pula \VUgras{batu pasir} yang dipotong
\VUgras{penggali batu} tua itu dengan \VUgras{beliungnya,} melainkan
hanya \VUgras{batu potong} sederhana untuk membangun rumah.

%%K t09-c1-14-1 p
14. --- Salah satu hiasan paling indah negeri itu adalah
\VUgras{pohon-pohon fir} \textsuperscript{(31)}. Di mana-mana — di
dasar \VUgras{jurang} \textsuperscript{(32)}, di tepi \VUgras{arus
deras} \textsuperscript{(33)}, di tepi \VUgras{jurang dalam}
\textsuperscript{(34)} atau \VUgras{tebing curam} — \VUgras{daun
jarumnya} yang halus memberi bayangan hijau.

%%K t09-tit-6 sub
\VUsaut{3.0mm}
\VUpk{10.2pt}{40}{Berjalan kaki.}
\VUsaut{3.0mm}

%%K t09-c1-15-1 p
15. --- Saya memanfaatkan liburan saya untuk \VUgras{berjalan-jalan}
jauh. \VUgras{Jalan-jalan} \textsuperscript{(35)} yang berkelok di
gunung memungkinkan menempuh, tanpa mempedulikan jaraknya, beberapa
\VUgras{kilometer} dan sering beberapa \VUgras{pal.}

%%K t09-c1-15-2 p
\VUgras{Batu-batu batas} \textsuperscript{(36)} dan \VUgras{tiang-tiang
penunjuk} \textsuperscript{(37)} menandai jalan-jalan itu, tempat orang
sering berjumpa dengan seorang \VUgras{orang gunung}
\textsuperscript{(38)} muda dengan \VUgras{buntil}
\textsuperscript{(40)} di sisinya, membawa seekor \VUgras{marmot}
\textsuperscript{(39)}, atau dengan seorang \VUgras{pengembara}
\textsuperscript{(41)} yang \VUgras{topi bulunya}
\textsuperscript{(42)} berbeda aneh dengan \VUgras{tudungnya}
\textsuperscript{(43)} tua yang koyak. Ia menuntun seekor
\VUgras{beruang berantai} \textsuperscript{(44)} yang sudah dilatih.

%%K t09-tit-7 sub
\VUpk{10.2pt}{40}{Mendaki gunung.}
\VUsaut{3.0mm}

%%K t09-c1-16-1 p
16. --- Hampir setiap hari saya pergi berlatih \VUgras{mendaki}
bersama seorang \VUgras{pemandu} \textsuperscript{(48)}, dan saya sangat
bangga ketika, dengan \VUgras{ransel} \textsuperscript{(47)} di
punggung dan « \VUgras{alpenstock} » di tangan, seperti seorang
\VUgras{wisatawan} \textsuperscript{(46)} sejati saya menyerbu
\VUgras{batu-batu karang} \textsuperscript{(45)}.

%%K t09-c1-17-1 p
17. --- Dari puncak gunung, penduduk dataran tampak tidak lebih besar
daripada semut ; sebuah \VUgras{kawanan domba} \textsuperscript{(50)}
yang merumput di \VUgras{padang rumput} \textsuperscript{(49)} tampak
seperti mainan anak-anak. Sebatang \VUgras{pohon pinus}
\textsuperscript{(51)} atau sebuah \VUgras{rumah kayu kecil}
\textsuperscript{(52)} tampak hampir hanya sebagai noda pada hijaunya
alam.

%%K t09-c1-18-1 p
18. --- Selama perjalanan itu kami sering berjumpa di \VUgras{jalan
setapak} \textsuperscript{(53)} dengan seorang \VUgras{pemburu chamois}
\textsuperscript{(54)} yang memikul di bahu binatang yang baru saja
dibunuhnya.

%%K t09-c1-19-1 p
19. --- Saya menyimpan di dalam buku tumbuhan saya sehelai dahan
\VUgras{paku-pakuan} \textsuperscript{(55)} yang sangat menarik dan
setangkai kecil \VUgras{bunga heather} \textsuperscript{(57)} merah
muda yang indah, yang saya petik di mulut sebuah \VUgras{gua}
\textsuperscript{(56)} kecil pada tamasya saya yang terakhir.

%%K t09-tit-8 sub
\VUsaut{3.0mm}
\VUcentre{12.0pt}{0}{\textit{Adegan Kedua.}}
\VUsaut{3.0mm}

%%K t09-tit-9 sub
\VUpk{11.4pt}{40}{Hutan. --- Perburuan.}
\VUsaut{3.0mm}

%%K t09-c2-01-1 p
1. --- Kemarin saya menghabiskan sehari yang nikmat di hutan.
Kerajinan yang merajalela di antara binatang dan \VUgras{serangga}
\textsuperscript{(5)} yang menghuninya sangat menarik perhatian saya.

%%K t09-c2-01-2 p
\VUgras{Laba-laba} \textsuperscript{(1)}, yang menenun \VUgras{jaringnya}
\textsuperscript{(2)} di antara dua dahan untuk menangkap
\VUgras{lalat} \textsuperscript{(3)} yang lewat ; \VUgras{siput}
\textsuperscript{(4)}, yang dengan susah payah memikul cangkangnya ;
kumbang bertanduk, yang memeriksa hasil tangkapannya ; semut yang
rajin, yang begitu bersemangat di dekat \VUgras{sarang semut}
\textsuperscript{(6)} — semua makhluk kecil itu pergi, datang dan
bekerja dengan kerajinan yang luar biasa.

%%K t09-c2-02-1 p
2. --- Seekor \VUgras{ular} \textsuperscript{(7)}, yang pada pandangan
pertama saya kira \VUgras{ular beludak} tetapi yang ternyata hanya
\VUgras{ular tidak berbisa} biasa, meringkuk di bawah sebatang pohon
dan sesekali menjulurkan kepala kecilnya yang ramping, yang selalu
tampak siap mematuk. Di depan saya seekor \VUgras{siput tanpa
cangkang} \textsuperscript{(8)} yang gemuk merayap dengan berat sesudah
melahap salah satu \VUgras{arbei hutan} \textsuperscript{(11)} yang
lezat, yang tumbuh melimpah di sana.

%%K t09-c2-03-1 p
3. --- Tanahnya berlapis \VUgras{bunga-bunga hutan} :
\VUgras{bunga primula} \textsuperscript{(9)}, \VUgras{bunga tapak dara}
\textsuperscript{(10)} dan banyak tanaman lain yang tidak saya kenal
bermekaran di antara \VUgras{rumpun-rumpun rumput}
\textsuperscript{(12)}.

%%K t09-c2-04-1 p
4. --- Di daerah itu \VUgras{trufel} hampir sama umumnya dengan
\VUgras{jamur} \textsuperscript{(14)}, dan seekor \VUgras{anak babi
hutan} \textsuperscript{(13)} milik seorang penjaga hutan sedang
menjelajah dengan rakus untuk melihat apakah ia akan berhasil
menemukan beberapa di antaranya.

%%K t09-tit-10 sub
\VUsaut{3.0mm}
\VUpk{10.2pt}{40}{Perburuan gelap.}
\VUsaut{3.0mm}

%%K t09-c2-05-1 p
5. --- Saya menyaksikan \VUgras{akhir} sebuah adegan yang sangat
menghibur saya. Dengan bantuan penciuman \VUgras{anjing bassetnya}
\textsuperscript{(15)}, \VUgras{penjaga hutan}
\textsuperscript{(16)} baru saja memergoki seorang \VUgras{pemburu
gelap} \textsuperscript{(22)} tepat pada saat orang itu bersiap membawa
pergi \VUgras{hewan buruan} \textsuperscript{(17)} yang telah
dibunuhnya. Karena lebih suka meninggalkan hasil buruannya daripada
tertangkap, pemburu gelap itu melarikan diri, dan seekor
\VUgras{kelinci hutan} \textsuperscript{(18)}, seekor \VUgras{rusa
betina} \textsuperscript{(19)}, seekor \VUgras{kijang}
\textsuperscript{(20)} dan seekor \VUgras{babi hutan}
\textsuperscript{(21)} tergeletak di rumput, menggembirakan penjaga
hutan.

%%K t09-c2-06-1 p
6. --- Beragamnya pohon yang ada di hutan itu mengherankan.
\VUgras{Pohon-pohon beech} \textsuperscript{(23)} dan \VUgras{pohon-
pohon birk} \textsuperscript{(26)}, \VUgras{pohon-pohon pinus,} yang
menghasilkan \VUgras{damar} \textsuperscript{(27)}, \VUgras{pohon-pohon
ek} \textsuperscript{(29)} dan masih banyak lagi yang lain
mencampurkan dahan-dahannya.

%%K t09-c2-06-2 p
\VUgras{Tanaman ivy} \textsuperscript{(24)}, yang melekat pada batang
pohon-pohon tinggi, mewarnai \VUgras{rimba} \textsuperscript{(25)}
dengan hijau yang berkilau, yang berpadu dengan menyenangkan dengan
warna \VUgras{lumut} \textsuperscript{(31)} yang lebih terang dan hijau
pucat, tempat \VUgras{burung pelatuk} \textsuperscript{(30)} mencari
serangga.

%%K t09-tit-11 sub
\VUsaut{3.0mm}
\VUpk{10.2pt}{40}{Pemandangan hutan.}
\VUsaut{3.0mm}

%%K t09-c2-07-1 p
7. --- Pohon-pohon ek tua memesona saya. \VUgras{Akar-akarnya}
\textsuperscript{(32)} yang besar dan berbelit, setengah keluar dari
tanah, memberi mereka rupa kekuatan dan tenaga yang megah, dan saya
bertanya-tanya bagaimana \VUgras{pemiliknya} \textsuperscript{(35)}
dapat merelakannya ditebang dan diserahkan kepada \VUgras{gergaji}
\textsuperscript{(34)} \VUgras{penebang} \textsuperscript{(33)}.
\VUgras{Batang-batang} \textsuperscript{(36)} yang malang itu, yang
dikupas \VUgras{kulit kayunya} \textsuperscript{(37)} atau dibelah oleh
\VUgras{kapak} \textsuperscript{(38)} \VUgras{buruh harian}
\textsuperscript{(39)}, membuat saya sedih.

%%K t09-c2-08-1 p
8. --- Di dekatnya, seorang \VUgras{pembuat arang}
\textsuperscript{(41)} tua telah menyiapkan dengan \VUgras{sekopnya}
sebuah \VUgras{tumpukan} \textsuperscript{(42)} arang, dan seorang
perempuan tua membawa \VUgras{seikat} \textsuperscript{(40)}
\VUgras{kayu mati} dari ranting-ranting pohon yang ditebang.

%%K t09-tit-12 sub
\VUsaut{3.0mm}
\VUpk{10.2pt}{40}{Perburuan.}
\VUsaut{3.0mm}

%%K t09-c2-09-1 p
9. --- Saya hendak pergi beristirahat beberapa saat di \VUgras{rumah
penjaga hutan} \textsuperscript{(46)} ; tetapi ketika saya sampai di
dekat \VUgras{danau kecil} \textsuperscript{(43)}, tempat seekor
\VUgras{angsa putih} \textsuperscript{(45)} yang indah berenang, saya
melihat bahwa \VUgras{banjir} \textsuperscript{(44)} yang baru saja
terjadi telah mengubah jalan itu menjadi \VUgras{rawa} yang
sesungguhnya. Saya terpaksa berbalik, dan sambil lewat di dekat
\VUgras{pohon cedar} \textsuperscript{(49)}, \VUgras{pohon-pohon
poplar} \textsuperscript{(48)} dan \VUgras{pohon elm}
\textsuperscript{(47)} yang berdiri di tepi hutan, saya melihat keluar
dari \VUgras{semak belukar} \textsuperscript{(54)} seekor \VUgras{rusa
jantan} \textsuperscript{(50)} yang sangat indah, yang dikejar
\VUgras{kawanan anjing pemburu} \textsuperscript{(51)} yang gigih, yang
juga dibakar semangatnya oleh \VUgras{terompet berburu}
\textsuperscript{(53)} \VUgras{para pemburu berkuda}
\textsuperscript{(52)}. Binatang yang malang itu benar-benar
kehabisan tenaga. Ia datang dan jatuh sekarat beberapa langkah dari
saya.

%%K t09-c2-10-1 p
10. --- Anak lelaki penjaga hutan, yang tertarik oleh keributan
\VUgras{perburuan berkuda} itu, keluar dari rumah, dan kami berdua
kembali ke hutan. Ia telah melihat seekor \VUgras{murai}
\textsuperscript{(56)} di dahan sebatang pohon ek tua. Ia mengira
sarangnya barangkali tersembunyi di \VUgras{dedaunan}
\textsuperscript{(58)}, dan ia hendak mengambil sarang itu.

%%K t09-c2-10-2 p
Setangkas \VUgras{tupai} \textsuperscript{(61)}, ia memanjat dari
\VUgras{dahan} \textsuperscript{(55)} ke dahan, tetapi ia tidak
menemukan apa-apa dan hanya berhasil menerbangkan seekor
\VUgras{burung gagak} \textsuperscript{(60)} tua yang hinggap di sebuah
\VUgras{ranting} \textsuperscript{(57)}.

\VUsaut{4.0mm}
\VUfilet{22mm}
